\documentclass[11pt]{article}
\usepackage[margin=1in]{geometry}
\usepackage{amsmath,amssymb,amsthm}
\usepackage{graphicx}
\usepackage{placeins}
\usepackage{booktabs}
\usepackage[expansion=false]{microtype}
\usepackage[hidelinks]{hyperref}
\setlength{\emergencystretch}{3em}

\theoremstyle{definition}
\newtheorem{proposition}{Proposition}

\title{\bf Topological Strain Memory on the $6N$ Skeleton:\\
the Effective Second-Order Markov Dynamics\\
of Prime Wing Transitions%
\thanks{Code and measured data: \url{https://github.com/Ruqing1963/6N-wing-transitions}}}
\author{Ruqing Chen\\
\small GUT Geoservice Inc., Montreal, Quebec, Canada\\
\small \texttt{ruqing@hotmail.com}}
\date{June 8, 2026}

\begin{document}
\maketitle

\begin{abstract}
Every prime $>3$ sits on a wing of the $6N$ lattice---right ($6N{+}1$) or left ($6N{-}1$)---so the
prime sequence is a binary string of wing symbols. Lemke Oliver and Soundararajan showed that
consecutive primes avoid repeating their residue class; here we ask a sharper question of the wing
string: \emph{how deep is its memory?} Using all $26{,}355{,}865$ primes $3<p\le5\times10^{8}$, we test
the Markov order of the chain. It is decisively \emph{not} first-order: the probability of staying on a
wing depends on the previous two primes, not one ($P(R\mid RR)=0.4252$ versus $P(R\mid LR)=0.4580$, a
gap of $0.033$ at $168\sigma$). The cleanest probe is the switch (step) sequence, which for a symmetric
first-order chain would be independent noise: instead it carries an autocorrelation of $-0.0328$ at lag
one (a switch is followed by a switch less often than chance---we name this a topological strain
memory), a weak $+0.0046$ echo at lag two, and \emph{nothing} measurable beyond. The incremental memory
contributed by each further prime halves per order ($0.033,0.016,0.009,\dots$), so the chain is
\emph{effectively}---not exactly---second-order: a dominant nearest-neighbour memory with a rapidly
decaying tail. We then show the wing memory and the well-known anti-correlation of consecutive prime
gaps are \emph{the same effect}: the switch indicator is identically the event ``gap $\not\equiv0
\pmod6$,'' and the step lag-one autocorrelation ($-0.0328$) equals the normalised-gap lag-one
correlation ($-0.0332$). All signals fade slowly with scale, as the Hardy--Littlewood framework
requires. This is a measurement, not a theorem: the structure lives inside the Hardy--Littlewood
$r$-tuple heuristic and the Oliver--Soundararajan bias, and what we add is the explicit memory-depth
characterisation and the gap identification.
\end{abstract}

\section{Introduction}
A prime $p>3$ is $\equiv1$ or $\equiv5\pmod6$; call these the right wing ($6N{+}1$) and the left wing
($6N{-}1$). Reading the primes in order turns them into a string $\dots\mathrm{RLRRLR}\dots$, and the
naive Cram\'er model says it should be a fair coin: each wing equally likely, transitions independent.
It is not. Lemke Oliver and Soundararajan \cite{OS} found that consecutive primes are biased
\emph{against} repeating their class modulo $q$; for $q=3$---which is exactly our wing label, since
$6N{+}1\equiv1$ and $6N{-}1\equiv2\pmod3$---this makes same-wing transitions rarer than switches.

The bias itself is by now expected. Our question is one level deeper and, to our knowledge, not usually
posed directly: \emph{what is the memory depth of the wing string?} It is sharpest to set three targets
against which to read the answer. The \emph{zeroth} target is the Cram\'er model: the wings are a fair
coin, switches unbiased and independent. The \emph{first} target is the naive reading of the
Oliver--Soundararajan bias: a switch bias exists, but lives entirely in the current step, so that the
chance of switching does not depend on how the current wing was reached---$P(\text{switch}\mid\text{previous step})$
is constant. The \emph{third} possibility, which the data select, is that the chain is genuinely
higher-order: the switch probability depends on the previous step, and the dependence decays across
further steps rather than terminating at once. The first target is broken decisively; the third is what
we measure, with the memory dominated by the second order and a geometrically decaying---not abruptly
truncated---tail. This is a sharp, finite, measurable question with an answer we did not know in
advance, and we answer it with the full prime census to $5\times10^{8}$.

We are explicit from the outset about what is and is not new. The existence of higher-order structure is
contained in the Hardy--Littlewood $r$-tuple heuristic \cite{HL}, which predicts the frequency of every
finite pattern of consecutive primes, and the Oliver--Soundararajan analysis \cite{OS} is built on it.
We prove no theorem and propose no new mechanism. What this paper contributes is the explicit
characterisation of the memory---its order, its sign, its decay rate---in the geometric language of the
$6N$ wings, and the identification of that memory with the anti-correlation of consecutive prime gaps.

\section{Setup: the wing string and its switch sequence}
Let $w_i\in\{\mathrm{R},\mathrm{L}\}$ be the wing of the $i$-th prime $>3$. The two wings occur almost
equally often (to $5\times10^{8}$, $13{,}177{,}500$ right against $13{,}178{,}365$ left; the residual
Chebyshev bias is below $10^{-4}$), so the chain is very nearly symmetric under $\mathrm{R}\!\leftrightarrow\!\mathrm{L}$.

This symmetry gives a clean test of Markov order. Define the \emph{switch} (step) sequence
$\sigma_i=\mathbf{1}[w_{i+1}\ne w_i]$. For a \emph{symmetric} first-order Markov chain the switch
probability is the same in both states, so $\sigma_i$ would be an independent Bernoulli sequence with
\emph{zero} autocorrelation at every lag. Any measured autocorrelation of $\sigma$ is therefore direct
evidence that the wing chain is \emph{not} first-order. We enumerate the primes by a sieve to
$N=5\times10^{8}$ and form $w$, $\sigma$, and the prime gaps $d_i=p_{i+1}-p_i$ directly; no modelling
enters.

\section{The chain is not first-order}
The marginal transition is $P(\text{same wing})=0.441$, $P(\text{switch})=0.559$: already the
$6\%$ Oliver--Soundararajan deficit against the fair-coin $0.5$. Conditioning on one more prime breaks
the first-order hypothesis outright (Table~\ref{tab:trip}): with the current wing held at R, the next
wing still depends on the one before,
\[
P(R\mid RR)=0.4252,\qquad P(R\mid LR)=0.4580,
\]
a difference of $0.0328$. With $5.8\times10^{6}$ samples in each cell the standard error is $\sim2\times
10^{-4}$, so the gap is significant at $168\sigma$. A first-order chain requires these to be equal; they
are not.

\begin{table}[t]
\centering\small
\begin{tabular}{lcc|lcc}
\toprule
pattern & count & $P(\text{next}=R)$ & pattern & count & $P(\text{next}=R)$\\
\midrule
RR & $5{,}843{,}102$ & $0.4252$ & LR & $7{,}334{,}397$ & $0.4580$\\
RL & $7{,}334{,}397$ & $0.5420$ & LL & $5{,}843{,}967$ & $0.5748$\\
\bottomrule
\end{tabular}
\caption{Second-order conditional probabilities. Holding the current wing fixed, the next wing still
depends on the previous one ($P(R\mid RR)\ne P(R\mid LR)$), so the chain is not first-order Markov.}
\label{tab:trip}
\end{table}

\section{The memory depth}
How far back does the dependence reach? The switch-sequence autocorrelation (Fig.~\ref{fig:fig1}(A))
answers cleanly, because under the first-order null it should vanish at all lags:
\[
\rho_1=-0.0328\ (168\sigma),\quad \rho_2=+0.0046\ (24\sigma),\quad \rho_3=-1\times10^{-5}\ (0\sigma),
\]
with $\rho_{\ge3}$ indistinguishable from zero. The memory is overwhelmingly nearest-neighbour: a switch
is \emph{less} likely to be followed by another switch ($P(\sigma{=}1\mid\sigma_{\text{prev}}{=}1)=0.544$
against $P(\sigma{=}1\mid\sigma_{\text{prev}}{=}0)=0.578$), as if a jump between wings relieves a strain
that then takes a step to rebuild---a kind of arithmetic inertia, the sequence resisting two abrupt
moves in a row. We use the image only as a name for the measured sign of $\rho_1$; it describes an
autocorrelation, not a force. The sign---anti-persistence of switches---is the dynamical face of
the Oliver--Soundararajan deficit.

``Effectively second-order'' must be stated carefully, and the full conditional table makes the
qualification honest. The incremental memory contributed by the oldest prime in a pattern---the average
of $|P(R\mid\mathrm{R}\cdots)-P(R\mid\mathrm{L}\cdots)|$ over the intervening symbols---does not cut off
abruptly; it \emph{halves} with each order (Fig.~\ref{fig:fig1}(B)):
\[
0.0328\ (m{=}2)\ \to\ 0.0159\ (m{=}3)\ \to\ 0.0092\ (m{=}4),
\]
ratios $0.48$ and $0.58$, descending toward the finite-sample floor by $m\approx5$. So the chain is not
exactly a second-order Markov chain; it is one with a dominant second-order term and a geometrically
decaying tail of weaker higher-order corrections. ``Effective second-order'' names this hierarchy, not a
hard truncation. The switch autocorrelation, which already collapses by lag three, and the
slower conditional decay are two views of the same short-ranged, rapidly fading memory.

\section{One effect, two projections: wings and gaps}
The wing memory is not separate from the known anti-correlation of consecutive prime gaps---it is the
same statistic. A switch occurs exactly when the gap is not a multiple of $6$: to $5\times10^{8}$ the
identity $\sigma_i=\mathbf{1}[\,6\nmid d_i\,]$ holds for every prime (agreement $1.0000$), because
crossing wings changes the residue mod $3$. Measuring both,
\[
\rho_1(\sigma)=-0.0328,\qquad \mathrm{corr}(g_i,g_{i+1})=-0.0332\quad(g_i=d_i/\log p_i),
\]
which agree to the third decimal. The conditional gap curve (Fig.~\ref{fig:fig2}(A)) shows the mechanism
in gap language: after a small gap the next gap runs above average, after a large gap below---primes
resist clustering, and a wing switch is the coarse, mod-$6$ shadow of that resistance. We state this as
an identification, not a metaphor: the second-order wing memory and the consecutive-gap
anti-correlation are one anti-clustering tendency seen through two resolutions.

\begin{figure}[t]
\centering
\includegraphics[width=\textwidth]{p34_fig1.pdf}
\caption{(A) Autocorrelation of the switch sequence: under a first-order chain it would be zero
everywhere. It is $-0.033$ at lag one ($168\sigma$), a weak $+0.005$ at lag two, and consistent with
zero beyond. (B) The incremental memory contributed by each further prime halves per order, descending
toward the finite-sample floor---a decaying tail, not an abrupt cutoff.}
\label{fig:fig1}
\end{figure}

\begin{figure}[t]
\centering
\includegraphics[width=\textwidth]{p34_fig2.pdf}
\caption{(A) Gap memory: the expected next normalised gap decreases with the current one (independence
would be flat at $1$). The wing switch lag-one autocorrelation and the normalised-gap lag-one
correlation coincide ($-0.0328$ vs $-0.0332$): one anti-clustering effect, two projections. (B) Both the
first-order ($\tfrac12-P(\text{same})$) and second-order ($|P(W|W)-P(W|S)|$) signals fade slowly as the
primes grow, as the Hardy--Littlewood framework requires.}
\label{fig:fig2}
\end{figure}

\FloatBarrier
\section{Scope and honest ledger}
The wing string of the primes is not a fair coin and not a first-order Markov chain. It carries a
nearest-neighbour memory---a switch suppresses the next switch by $0.033$---with a weak second echo and
a geometrically decaying higher-order tail, so it is effectively, not exactly, second-order. That memory
is the same object as the anti-correlation of consecutive gaps, the switch being identically the event
that the gap is not divisible by $6$. Every signal fades slowly toward the fair-coin limit as the primes
grow (Fig.~\ref{fig:fig2}(B)): $P(\text{same})$ rises from $0.404$ below $10^{5}$ to $0.446$ near
$5\times10^{8}$, and the second-order term shrinks in step.

We are explicit about the boundary. The anti-repetition bias is Lemke Oliver and Soundararajan's
\cite{OS}; the heuristic that predicts every consecutive-pattern frequency, and hence all of the
structure measured here, is Hardy and Littlewood's \cite{HL}; the slow fade with scale is the expected
behaviour of these conjectural densities. We prove nothing and propose no mechanism beyond these. What
this paper adds is the explicit memory-depth characterisation of the wing chain---its order, the sign
and magnitude of each term, and the geometric decay---together with the identification of that memory
with the consecutive-gap anti-correlation, all in the $6N$ geometric language and measured against the
full prime census to $5\times10^{8}$. The words ``strain'' and ``anti-clustering'' name measured
correlations; they are not claims about a physical force. Everything here equals a direct count over the
primes and takes the Hardy--Littlewood heuristic as input.

\paragraph{Data and code availability.} The sieve to $5\times10^{8}$, the wing and switch sequences, the
order-$1$ to order-$4$ transition tables, the autocorrelation and memory-depth measurements, the
gap identification, and the figure scripts are openly available at
\url{https://github.com/Ruqing1963/6N-wing-transitions}.

\begin{thebibliography}{9}
\bibitem{HL} G.~H.~Hardy, J.~E.~Littlewood, \emph{Some problems of `Partitio Numerorum'; III}, Acta Math.\ \textbf{44} (1923), 1--70.
\bibitem{OS} R.~J.~Lemke Oliver, K.~Soundararajan, \emph{Unexpected biases in the distribution of consecutive primes}, Proc.\ Natl.\ Acad.\ Sci.\ USA \textbf{113} (2016), E4446--E4454.
\bibitem{PXXIX} R.~Chen, \emph{Twin counts in short intervals: singular-series rigidity and the density-gradient variance} (Part XXIX), Zenodo, doi:10.5281/zenodo.20597725 (2026).
\bibitem{PXXXIII} R.~Chen, \emph{The Fermat-type dodge ladder: primes $n^{2^k}+1$ on the 6N skeleton} (Part XXXIII), Zenodo, doi:10.5281/zenodo.20602622 (2026).
\bibitem{review} R.~Chen, \emph{Arithmetic geodynamics on the 6N skeleton: a review of Parts I--XIX}, Zenodo, doi:10.5281/zenodo.20585301 (2026).
\end{thebibliography}

\end{document}
